\chapter{Digitisation with Shannon}
\label{ch:digitisation}

\chaptergoals{9 of 11 in Part~II.}
{\cref{ch:aegir}: \file{muon\_gun.root}.}
{digitise simulated hits, explain what the hit-stream output means for an analysis, and test the smearing.}

\section{From energy deposits to detector hits}

Geant4 tells us where charged particles deposited energy. A detector reports
something else: which channel fired, when, and with what signal, after
thresholds, noise, dead time, and the merging of nearby deposits into one
channel. \emph{Digitisation} is the step that turns the first into the second,
and \repo{Shannon} is the package that does it for SHiP. Its output provides the hit-stream interface being developed for reconstruction;
a realistic electronics response and a real-data adapter are separate tasks.

\pkg{Shannon} is a set of Phlex plugins, like \pkg{aegir}:

\begin{center}
\small
\begin{tabularx}{\linewidth}{l l L}
\toprule
Plugin & Kind & Purpose \\
\midrule
\texttt{read\_sim\_file}            & source   & reads \texttt{sim\_hits} and \texttt{sim\_particles} from \pkg{aegir} output \\
\texttt{digitise\_hits}             & transform& one digitiser per detector, chosen by \field{detector\_id} \\
\texttt{digitised\_output\_module}  & observer & one RNTuple per hit collection, validation histograms \\
\bottomrule
\end{tabularx}
\end{center}

\section{What the digitisers do today}

Every detector has its own digitiser class in \file{src/detectors/}, so that
each detector group can develop its model independently. At the time of
writing all five share a placeholder model:

\begin{itemize}
  \item each \texttt{SimHit} becomes exactly one detector hit;
  \item $x$ and $y$ are smeared with a Gaussian of $\sigma = 0.1$\,mm
        (\file{gaussian\_smearer.hpp}); $z$, the energy deposit, the momentum
        and the truth links (\field{track\_id}, \field{pdg\_code}) are copied;
  \item a time offset is added: \pkg{Shannon} places each simulated event at a
        random time within the fraction of a spill that the simulated protons
        on target correspond to (the \texttt{pot} parameter; a SPS spill is
        1.2\,s for $4\times10^{13}$ protons). Our artificial muons do not
        come from a number of protons, so here \texttt{pot} only sets the
        window: with the 1000 protons of \cref{ch:production} it is 30\,ns,
        the time those protons take to arrive. For a replay of a real
        production, \texttt{pot} is the number of protons of stage~1.
\end{itemize}

There are no thresholds, no channel merging, no noise and no inefficiency yet.
The structure is the point of this chapter, not the realism.

\begin{note}
Random numbers come from \pkg{ship-random}: two \texttt{PhiloxRng} streams per
spill (one for the time offset, one for the smearing), keyed on the spill
number. Digitisation is therefore reproducible for a given \texttt{seed},
whatever the number of threads for a fixed input entry order. Reordering
the simulation entries changes the event-to-random-stream assignment.
\end{note}

\section{Running it}

The workflow is \file{code/workflows/ft\_digi.jsonnet}:

\codefile[linerange={19-66},basicstyle=\ttfamily\footnotesize]{jsonnet}{code/workflows/ft_digi.jsonnet}

\begin{exercise}{Digitise the muons}
In the \pkg{Shannon} environment, from your working directory, with the 100
entries of \file{muon\_gun.root} (\cref{ex:muongun}) and the 1000 protons of
\cref{ex:production}:
\begin{shell}
cd "$SHIP_TUTORIAL/Shannon"
pixi cvmfs
cd "$SHIP_TUTORIAL/work"
phlex -j 4 -c <(jsonnet --tla-str infile=muon_gun.root \
    --tla-code entries=100 --tla-code pot=1000 \
    $CODE/workflows/ft_digi.jsonnet)
\end{shell}
\end{exercise}

\begin{checkpoint}
\file{ft\_digi.root} and \file{ft\_digi\_validation.root} exist. In
Python, \texttt{ROOT.TFile("ft\_digi.root").ls()} lists seven RNTuples, and
\texttt{h\_straw\_tubes\_multiplicity} in the validation file shows the number
of straw hits per spill.
\end{checkpoint}

\begin{pitfall}
``No digitiser registered for detector ID $n$'' means the input contains hits
of a detector Shannon does not know, for example the scoring-plane hits
(\field{detector\_id}~100) of a stage-1 file: digitise \file{muon\_gun.root},
not the production.
\end{pitfall}

\section{A stream of hits, not a list of events}
\label{sec:digilayout}

\pkg{Shannon} writes every collection as its own RNTuple with one entry per
hit (\cref{sec:filelayout}). This is deliberate. SHiP reads out its detectors
continuously, without a trigger: what the experiment records is a stream of
time-stamped hits per detector, and grouping hits that belong together into
candidate events, by time, is a task of reconstruction. The hit-stream interface is intended for both simulated and real inputs;
that design does not make today's placeholder response equivalent to real data.

For an analysis of \emph{simulated} events this has three consequences:

\begin{enumerate}
  \item \textbf{There is no event number.} A hit does not record which
        simulated event it came from, and the writers fill from several
        threads, so the order of entries does not follow the events either.
  \item \textbf{The generator record and the event weight are not copied.}
        \texttt{mc\_particles} and \texttt{event\_header} stay in the
        simulation file.
  \item \textbf{Truth links are kept per hit:} \field{track\_id} and
        \field{pdg\_code} survive in \texttt{RecHit}, but track IDs are only
        unique within one simulated event.
\end{enumerate}

Inclusive distributions over all hits, such as positions, times, energy
deposits and resolutions, come from the digitised file. Questions about one
simulated event, such as ``this muon crossed four straw stations'' or
``the SBT saw it'', are answered from the simulation file.
\Cref{ch:analysis} is organised accordingly.

\begin{organisernote}
The hit-stream layout is intended (see the header comment of
\file{Trout/src/io/read\_digitised\_hits.cpp}). What is missing for Monte
Carlo studies is a way back from a digitised hit to its simulated event. Even
with \texttt{pot} set to the protons of stage~1, all events share a window of
$1.2\,\mathrm{s}\times \texttt{POT}/(4\times10^{13})$, a few tens of
nanoseconds for a thousand protons, and their hits overlap in time. A
simulated-event index carried in the hit (or a truth-link collection) would
let the analysis in \cref{ch:analysis} use digitised hits event by event.
\pkg{Shannon}'s pull request~\#26 goes in that direction. Worth raising with the \pkg{Shannon} developers; it is also a
good topic for the closing discussion.
\end{organisernote}

\begin{exercise}{Check the smearing}
Since \pkg{Shannon} copies $z$, the energy deposit and the truth links, many
digitised hits can be paired by a composite key. Repeated keys are ambiguous
because there is no event ID; the script excludes them. The script
\file{code/analysis/digi\_residuals.py} does this for one collection and
histograms the residual $x_\mathrm{digi} - x_\mathrm{sim}$:
\begin{shell}
python $CODE/analysis/digi_residuals.py ft_digi.root \
    --collection straw_tubes_hits
\end{shell}
Fit the residual distribution. Does the width match the smearing in
\file{gaussian\_smearer.hpp}? Repeat for \texttt{sbt\_hits}, where most of
the hits of the muon background are.
\end{exercise}

\begin{solution}
With sufficient uniquely matched hits, both fitted widths should be
compatible with 0.1\,mm, because all digitisers use the same default
\texttt{GaussianSmearer}. Some hits may stay unpaired when two
simulated hits share the same key; the script reports how many.
\end{solution}

\begin{exercise}[colback=white]{Optional: make the straw model your own}
Real straws measure the distance of closest approach to the wire with a
resolution of roughly 100--150\,\textmu m, perpendicular to the wire only.
In your \pkg{Shannon} checkout, give the straw digitiser its own resolution:
in \file{src/detectors/straw\_tubes.hpp} construct the smearer as
\texttt{GaussianSmearer smearer\_\{0.12, 0.12\};}, then
\begin{shell}
pixi cvmfs build
\end{shell}
rerun the digitisation and the residual script, and check that only the
straws changed. What would a realistic model need to know that a
\texttt{SimHit} does not provide directly?
\end{exercise}

\begin{solution}
The straw (wire) identity and orientation, which the hit's
\field{geometry\_node\_id} or position must be mapped to through the geometry
service; the drift-time relation; and the readout's time window. This is the
kind of question that connects the digitisation to the geometry service and
is a good lead-in to the reconstruction discussion.
\end{solution}
